\section{Results}\label{sec:results}

The structural estimates compress into four objects per class. The total price effect $p_{S_3} - p_{S_1}$ is $+0.24$ of the reference price in non-pharma and $+0.31$ in pharma, with 95 percent bootstrap CIs of $[0.19, 0.29]$ and $[0.25, 0.36]$. The intensive-margin share $(p_{S_2} - p_{S_1})/(\text{total})$ is 76 percent in non-pharma and 75 percent in pharma; the entry-margin share is the complement. A Turnbull NPMLE ratifies these numbers within five percent. And a counterfactual 10 percent SME price preference---developed in Section~\ref{sec:pareto}---delivers essentially zero price cost while preserving SME win-rate gains of the same order as the set-aside. The reduced-form precursor to this paper reported 10--11 percent price effects on winning prices; the structural estimates confirm the sign, materially raise the magnitude once auction-level heterogeneity and the pool-restriction mechanism are made explicit, and add the counterfactual instrument that the reduced form could not.

\subsection{Entry response}

Table~\ref{tab:v3_entry_rates} reports the panel. Pre-policy Preg\~ao has 0.86 SMEs and 2.46 non-SMEs per auction in non-pharma, 0.44 and 2.56 in pharma. Post-policy the counts rotate: 1.74 SMEs and 1.40 non-SMEs in non-pharma, 1.09 and 1.62 in pharma. SMEs roughly double in both classes. Non-SMEs drop by 43 percent in non-pharma and 37 percent in pharma. The doubling is the extensive-margin force that, in the BNE decomposition below, partially but incompletely offsets the intensive-margin price rise.

\subsection{BNE price decomposition}

Table~\ref{tab:v3_bne_decomp} gives the point estimates under the all-bidders regime and endogenous entry---the v3 baseline specification. In non-pharma, the simulated open-regime price is $p_{S_1} = 0.76$ of the reference price. Holding the SME pool at its pre-period size and restricting to SME-only raises the simulated price to $p_{S_2} = 1.11$---an intensive shift of +0.355. Letting SMEs enter at their post-policy rate pulls the simulated price down to $p_{S_3} = 1.01$---an entry-margin pullback of $-0.106$. The net effect is $p_{S_3} - p_{S_1} = +0.248$.

The pharma numbers follow the same pattern. $p_{S_1} = 0.66$, $p_{S_2} = 1.17$, $p_{S_3} = 1.00$, total $+0.341$. The intensive shift is $+0.510$; the entry offset is $-0.169$. The larger magnitudes in pharma reflect the heavier auction-level heterogeneity in that class (intraclass correlation 0.62 versus 0.47) and the thinner SME baseline pool ($N^{\text{SME}}_{\text{Pre}} = 0.44$ versus $0.86$ in non-pharma), both of which amplify the intensive component.

Dividing the intensive effect by the sum of absolute intensive and entry effects gives the intensive share: 77 percent in non-pharma, 75 percent in pharma. The common reading is that the set-aside operates roughly three quarters through bidding behavior and one quarter through entry behavior---the sign and order match the reduced-form intuition of the auction-preferences literature since \citet{marion2007}, but the quantitative split was not available in that literature.

\subsection{Sources of the departure from earlier estimates}

Table~\ref{tab:v3_decomp_grid} reports the $2 \times 2$ grid in which each axis is a methodological choice: cost distributions raw (losers-only, v2-style) versus UH-clean (v3), and entry fixed at the pre-period pool (v2-style counterfactual) versus endogenous (v3). The v3 headline sits in the bottom-right cell. The v2-style proxy sits in the top-left. The gap between them is 72 percent in non-pharma and 68 percent in pharma---that is, the earlier specification overstated the total effect by these amounts.

UH correction accounts for more than half of the gap: it drops the pharma intensive share from 87 percent (fixed-pool, raw) to 70 percent (fixed-pool, UH-clean), and similarly in non-pharma. Endogenous entry accounts for the rest. The two methodological additions operate in the same direction---both pull the headline down---and together bring the BNE price effect into the range that the reduced-form DiD in the predecessor paper supported (10--11 percent of absolute price, which corresponds to roughly 23--33 percent of reference, consistent with the v3 numbers).

\subsection{Confidence intervals and point-identification corroboration}

A cluster bootstrap at the auction level with B=500 replicates (Table~\ref{tab:v3_bootstrap_ci}) gives the following 95 percent CIs on $p_{S_3} - p_{S_1}$:
\begin{center}
\begin{tabular}{lcc}
\toprule
Class      & Mean     & 95\% CI \\
\midrule
Non-pharma & +0.236   & $[0.186, 0.289]$ \\
Pharma     & +0.305   & $[0.247, 0.364]$ \\
\bottomrule
\end{tabular}
\end{center}
The intensive-margin 95\% CI is $[65\%, 88\%]$ in non-pharma and $[63\%, 83\%]$ in pharma; neither CI on total or on share crosses zero. The CI widths of about ten percentage points on the total and fifteen to twenty on the share reflect sampling variation at the auction level---the cluster structure respects the within-auction correlation induced by $a_t$.

The Turnbull NPMLE---which does not rely on Assumption~\ref{a:uh}---delivers $+0.242$ (non-pharma) and $+0.338$ (pharma), well within each CI (Table~\ref{tab:v3_sensitivity_fc}). Turnbull intensive shares are 74 percent in non-pharma and 81 percent in pharma; the latter is 6 percentage points above the all-bidders point, if anything strengthening the intensive-margin interpretation in pharma. The losers-only estimator, reported as a v2-proxy, gives $+0.253$ and $+0.343$, moving in the direction expected from its upward tail bias. The three regimes bracket a narrow range, and the main text reports the all-bidders number because of its intermediate bias direction and its computational simplicity.

\subsection{Entry cost calibration}

Table~\ref{tab:v3_entry_cost} reports the zero-profit implied entry cost per bidder. In non-pharma, $\kappa^{\text{SME}} =$ R\$0.55 and $\kappa^{\neg} =$ R\$2.46---roughly four percent and eighteen percent of the pharma-class-median reference price of R\$13.98. In pharma, $\kappa^{\text{SME}} =$ R\$0.11 and $\kappa^{\neg} =$ R\$0.51, again in the low-single-digit percent of the reference. The five-to-one ratio between non-SME and SME entry costs is consistent across classes and plausible given that non-SMEs typically operate at larger scale with more complex documentation requirements, both of which raise the per-bid overhead of responding to a BEC tender. The magnitudes are reported in ref-price units in the table and translated to reais at the pharma-class-median reference price; neither the ratio nor the policy implications depend on the translation choice.

\subsection{A Pareto-Improvement Counterfactual}\label{sec:pareto}

The natural alternative to a full set-aside is a price preference: SME bids are scored with a discount for the purpose of winner selection, but the government pays the actual bid. The US Small Business Administration has used this instrument since the 1970s; Brazilian federal procurement law allows it but leaves it to administrative discretion. Within the model, a 10 percent price preference delivers
\begin{equation}
p - p_{S_1} \;=\; +0.015 \text{ (non-pharma)},\qquad -0.004 \text{ (pharma)},
\end{equation}
against $p_{S_3} - p_{S_1} = +0.24$ and $+0.31$ under the full set-aside. The price cost to the government is \emph{essentially zero}. The SME win-rate gain is of the same order as under the set-aside---SMEs are selected as the preferred bidder in a comparable share of auctions---but the payment rule does not pull the equilibrium price up with them. Table~\ref{tab:v3_apv} reports the full comparison across four policy variants.

The mechanism is straightforward. Under the set-aside, non-SMEs are absent from the winning-price formation and the second-lowest cost is drawn from a shallower distribution. Under the price preference, non-SMEs remain active; they continue to anchor $c_{(2)}$ even when a preferred SME wins, because an SME wins only when her discounted cost is below every non-SME's cost. The price distribution is therefore indistinguishable from the open-regime distribution at the second order, while SMEs capture the allocation slot whenever they are within 10 percent of the non-SME minimum.

I report the Pareto-improvement claim with two caveats. First, the calculation is Vickrey-equivalent: in a first-price sealed-bid auction with the same preference, equilibrium bidding would adjust in ways \citet{maskinriley2000} characterize---SMEs bidding more aggressively on the protected margin, non-SMEs less aggressively on the competitive margin. The direction of the adjustment typically \emph{reduces} SME gains and \emph{increases} the government's savings, so the Pareto-improvement claim carries through in sign; only the magnitude depends on the strategic response, which this paper does not solve. Second, V3 is not bootstrapped for a confidence interval. The point estimate is stable across 2{,}000 Monte Carlo draws, but V3's inferential claim is weaker than the main decomposition's. I treat V3 as a within-model policy counterfactual rather than a fully inferred estimate.

What V3 adds to the paper is directional rather than quantitative. The structural decomposition (intensive vs.\ entry) is the technical core; the welfare cost in Section~\ref{sec:welfare} is the accounting; V3 is the policy corollary. The number 2.4--3.0---the welfare weight on SME surplus required to prefer the set-aside over the 10 percent preference---is developed in Section~\ref{sec:discussion}. Outside an implausible range of distributional preferences, the current rule is welfare-dominated by an instrument the law already allows.
